\section{What would count against the framework?}

A useful framework must identify evidence that would make it unnecessary or wrong.

The present proposal would be weakened if:
\begin{enumerate}[label=(\arabic*)]
\item current state and simple empirical trend consistently predict future failure as well as explicit measures of corrective capacity and renewal;
\item purported maintenance processes can be removed without affecting future viability under the disturbances they are claimed to handle;
\item independent failure topology adds no predictive information beyond nominal redundancy;
\item systems remain robust despite persistent collapse of all identifiable capacity-renewal pathways without drawing on external support or inherited reserves;
\item measured deterioration in function--proxy coupling fails to affect allocation or future viability;
\item vector-valued maintenance variables can consistently be replaced by a naturally determined scalar objective without loss of prediction;
\item the same observations are explained more parsimoniously by established viability, reliability or control models without the distributed-maintenance decomposition adding discriminating variables.
\end{enumerate}

The seventh possibility is especially important. The framework should not survive merely by redescribing known effects in new vocabulary.

\section{Discussion}

The central result of the modelling programme is not a new universal dynamical law. It is a sequence of restrictions on what can safely be inferred from what we observe.

A currently successful process need not supply enough control. A persistent organization need not currently reproduce the conditions responsible for its persistence. A present participant need not contribute positively after network response. A healthy present state need not imply spare corrective capacity. A reinforced process need not still perform the function for which reinforcement is useful. And contribution to one requirement need not determine contribution to another.

These restrictions shift the natural unit of analysis. Instead of asking whether a node, species, validator, institution or behaviour is \`\`good for the system,'' we ask:
\[
\boxed{
\begin{array}{l}
\text{Which viability condition is being considered?}\\
\text{What disturbance threatens it?}\\
\text{Which processes can correct that disturbance?}\\
\text{What failure domains do those processes share?}\\
\text{What maintains their corrective capacity?}\\
\text{Which signal allocates renewal to them?}\\
\text{How tightly does that signal remain coupled to function?}\\
\text{What other viability margins are changed at the same time?}
\end{array}
}
\]

This is a more demanding description than persistence alone. It is also experimentally richer.

\subsection{From control to maintenance}

Control usually emphasizes how an action changes a target variable. Maintenance adds time.

The controller itself must remain available. Its information source must remain reliable. Its resources must be replenished. Its failure modes must remain sufficiently independent. And its architecture may eventually need revision.

Thus
\[
\boxed{
\text{maintenance}
=
\text{control extended through the conditions required for future control}.
}
\]

This definition is deliberately functional. It does not require the maintaining processes to know that they are maintaining anything.

\subsection{From objects to processes}

The framework also shifts emphasis from enduring objects to recurrent process organization.

Individual components may be replaced while the relations producing future components persist. A maintenance architecture can therefore survive component turnover. Conversely, the same components can remain physically present after the relations that made them collectively functional have deteriorated.

What continues is not necessarily material identity. It is the organized capacity to regenerate viability-relevant relations.

\subsection{From a tree of causes to a web of maintenance}

Distributed maintenance is rarely a simple hierarchy. A process may maintain a resource that maintains another process that alters the first process's environment. Correction can be heterarchical. Signals can be generated by the same activity they regulate. Multiple feedback and return paths can overlap.

The appropriate representation is consequently closer to a web than a chain.

This does not mean that every causal network is self-maintaining. The empirical task is precisely to identify which edges close the relevant viability loops and which are incidental.

\section{Conclusion}

The existence of an organization tells us that some historical sequence produced its present state. It does not tell us whether the processes operating now are sufficient to preserve that organization, its functions, or its ability to respond to future disturbance.

To study that question we must distinguish at least four things:
\[
\boxed{
\text{present state},
\quad
\text{viability margin},
\quad
\text{corrective capacity},
\quad
\text{renewal of corrective capacity}.
}
\]
Distributed systems add two further complications:
\[
\boxed{
\text{independent access to correction}
\quad\text{and}\quad
\text{fidelity of the signal allocating renewal}.
}
\]

Taken together, these distinctions motivate the concept of \emph{distributed maintenance}: local processes alter the future viability of shared enabling conditions while themselves depending on information, resources, structure and other processes for continued corrective capacity.

No individual element of this account is claimed as unprecedented. Viability theory supplies the mathematics of constrained continuation. Organizational autonomy supplies a developed theory of mutually maintained constraints. Biological regulation supplies the concept of second-order internal control. Reliability engineering supplies common-mode failure. Adaptive-network theory supplies physical examples of use-dependent reinforcement.

The proposed contribution is their diagnostic composition around a common question:

\begin{quote}
\textbf{What must continue to work now for this system to remain capable of keeping itself within the conditions we have declared necessary for its future continuation?}
\end{quote}

Answering that question requires more than observing that the system still exists. It requires measuring the maintenance of the ability to continue.

\section*{Code, models and provenance}

The mathematical and computational exercises motivating this synthesis are maintained in the \emph{Distributed Commons Control} repository. Several boundary conditions used here were first developed in separate manuscripts in the \emph{Evolution by Emergence} repository. This paper intentionally distinguishes those prior results from the synthesis proposed here. Lean verification in those repositories establishes consequences of declared mathematical assumptions; it does not validate the empirical interpretation of the models.
